\documentclass[11pt]{article}
\usepackage[margin=1in]{geometry}
\usepackage{graphicx}
\usepackage{booktabs}
\usepackage{array}
\usepackage{enumitem}
\usepackage{caption}

% Simple line for fill-in blanks
\newcommand{\fillin}{\rule{2.5em}{0.4pt}}

\title{MIPS Datapath Lab Report Template}
\author{Your Name Here}
\date{CS161 -- Lab 2 and Lab 3}

\begin{document}
\maketitle

\section*{How to use this template}
\begin{itemize}[leftmargin=*]
  \item Replace placeholders (\texttt{\textbackslash fillin} and boxed notes) with your answers, figures, and screenshots.
  \item Keep all screenshots readable. If needed, enlarge critical probes or labels before capturing.
  \item Remember to rename the Logisim file to \texttt{lastname-firstname-mips-datapath-lab.circ} and this PDF to \texttt{lastname-firstname-mips-datapath-lab.pdf} before submitting.
  \item When you add graphics, swap each boxed placeholder with \verb|\includegraphics[width=\linewidth]{path/to/figure}| and update captions.
\end{itemize}

\section{ToDo \#1: Type Decode Unit}
\begin{itemize}[leftmargin=*]
  \item Briefly describe how you detected each instruction type (comparators, constants used, wiring choices).
  \item Verify each opcode below activates the intended output.
\end{itemize}

\begin{table}[h]
  \centering
  \caption{Opcode patterns for instruction types (given in handout).}
  \begin{tabular}{lcccc}
    \toprule
    Input (opcode bit) & R-format & lw & sw & beq \\
    \midrule
    Op5 & 0 & 1 & 1 & 0 \\
    Op4 & 0 & 0 & 0 & 0 \\
    Op3 & 0 & 0 & 1 & 0 \\
    Op2 & 0 & 0 & 0 & 1 \\
    Op1 & 0 & 1 & 1 & 0 \\
    Op0 & 0 & 1 & 1 & 0 \\
    \bottomrule
  \end{tabular}
\end{table}

\begin{figure}[h]
  \centering
  \fbox{\parbox{0.9\linewidth}{Screenshot placeholder: completed Type Decode circuit.}}
  \caption{Type Decode unit after adding opcode comparators.}
\end{figure}

\section{ToDo \#2: Control Decode Truth Table}
Fill in the control signals for each instruction type (use X for don't care).

\begin{table}[h]
  \centering
  \caption{Control-decode truth table to complete.}
  \begin{tabular}{lcccc}
    \toprule
    Output & R-format & lw & sw & beq \\
    \midrule
    RegDst   & 1     & 0     & \textit{X} & \textit{X} \\
    ALUSrc   & \fillin & \fillin & \fillin & \fillin \\
    MemtoReg & \fillin & \fillin & \fillin & \fillin \\
    RegWrite & \fillin & \fillin & \fillin & \fillin \\
    MemRead  & \fillin & \fillin & \fillin & \fillin \\
    MemWrite & \fillin & \fillin & \fillin & \fillin \\
    Branch   & \fillin & \fillin & \fillin & \fillin \\
    ALUOp1   & 1     & 0     & 0     & 0 \\
    ALUOp0   & 0     & 0     & 0     & 1 \\
    \bottomrule
  \end{tabular}
\end{table}

\section{ToDo \#3: Control Decode Implementation}
\begin{itemize}[leftmargin=*]
  \item Summarize how you generated the logic (manual derivation vs. Logisim Combinational Analysis).
  \item Note any simplifications or don't-care usage.
\end{itemize}

\begin{figure}[h]
  \centering
  \fbox{\parbox{0.9\linewidth}{Screenshot placeholder: completed Control Decode circuit.}}
  \caption{Control Decode unit with all outputs wired.}
\end{figure}

\section{ToDo \#4: Control Logic Wiring}
\begin{itemize}[leftmargin=*]
  \item Confirm connections from Type Decode outputs to Control Decode inputs, and from Control Decode outputs to the datapath signals.
  \item Mention any probe values you used to validate the chain end-to-end.
\end{itemize}

\begin{figure}[h]
  \centering
  \fbox{\parbox{0.9\linewidth}{Screenshot placeholder: fully wired Control block.}}
  \caption{Completed control logic wiring.}
\end{figure}

\section{ToDo \#5: ALU Control Unit}
\begin{itemize}[leftmargin=*]
  \item Describe the conditions you used to drive the priority encoder inputs (per the ALUOp/Func table below).
  \item Call out how you handled the cases that map to operations 0010, 0110, 0000, 0001, and 0111.
\end{itemize}

\begin{table}[h]
  \centering
  \caption{ALU control mapping (provided in handout).}
  \begin{tabular}{ccccccccc}
    \toprule
    ALUOp1 & ALUOp0 & F5 & F4 & F3 & F2 & F1 & F0 & Operation \\
    \midrule
    0 & 0 & X & X & X & X & X & X & 0010 \\
    0 & 1 & X & X & X & X & X & X & 0110 \\
    1 & 0 & X & X & 0 & 0 & 0 & 0 & 0010 \\
    1 & 0 & X & X & 0 & 0 & 1 & 0 & 0110 \\
    1 & 0 & X & X & 0 & 1 & 0 & 0 & 0000 \\
    1 & 0 & X & X & 0 & 1 & 0 & 1 & 0001 \\
    1 & 0 & X & X & 1 & 0 & 1 & 0 & 0111 \\
    \bottomrule
  \end{tabular}
\end{table}

\begin{figure}[h]
  \centering
  \fbox{\parbox{0.9\linewidth}{Screenshot placeholder: completed ALU Control circuit.}}
  \caption{ALU Control with priority encoder and mux logic filled in.}
\end{figure}

\section{ToDo \#6: Datapath Wiring}
\begin{itemize}[leftmargin=*]
  \item Note key connections (e.g., register destination MUX, ALU source MUX, MemtoReg MUX, PC path).
  \item Mention any probes you added or relied on for verification while stepping through instructions.
\end{itemize}

\begin{figure}[h]
  \centering
  \fbox{\parbox{0.9\linewidth}{Screenshot placeholder: fully wired datapath.}}
  \caption{Completed processor wiring (datapath view).}
\end{figure}

\section{ToDo \#7: Program Execution Results}
\begin{itemize}[leftmargin=*]
  \item Load \texttt{1234.mem} into data memory and \texttt{test-code.mem} into instruction memory, then tick through execution.
  \item Record final register values for \$t0--\$t6 after the program halts.
\end{itemize}

\begin{table}[h]
  \centering
  \caption{Final register file contents to report.}
  \begin{tabular}{lccccccc}
    \toprule
    Register & $t0$ & $t1$ & $t2$ & $t3$ & $t4$ & $t5$ & $t6$ \\
    \midrule
    Value & \fillin & \fillin & \fillin & \fillin & \fillin & \fillin & \fillin \\
    \bottomrule
  \end{tabular}
\end{table}

\begin{figure}[h]
  \centering
  \fbox{\parbox{0.9\linewidth}{Screenshot placeholder: register file after program completes.}}
  \caption{Register file state after executing all instructions.}
\end{figure}

\subsection*{Optional trace (helpful for debugging)}
\begin{table}[h]
  \centering
  \caption{Per-instruction register write trace.}
  \begin{tabular}{llll}
    \toprule
    Cycle & Instruction & Destination register & Value written \\
    \midrule
    1 & \fillin & \fillin & \fillin \\
    2 & \fillin & \fillin & \fillin \\
    3 & \fillin & \fillin & \fillin \\
    4 & \fillin & \fillin & \fillin \\
    5 & \fillin & \fillin & \fillin \\
    6 & \fillin & \fillin & \fillin \\
    7 & \fillin & \fillin & \fillin \\
    \bottomrule
  \end{tabular}
\end{table}

\section*{Debugging notes (use as needed)}
\begin{itemize}[leftmargin=*]
  \item Record any wiring pitfalls, mux select conventions, or probe readings that were important for fixes.
  \item Note anything you would check first if you had to re-debug the design.
\end{itemize}

\end{document}
